\section{Can the probabilities be trusted?}
\label{sec:calib}
\textbf{Code:} \code{calibration.py}, \code{experiments/task3\_calibration.py}.

Section~\ref{sec:sweep} found logistic regression cheaper than the Gaussian models with every rule.
To see why, we looked at the probabilities themselves. Accuracy only asks whether $\post$ is on the
right side of $0.5$. Our rule compares $\post$ with
$0.3$ and $0.97$, so it needs $\post$ to be \emph{calibrated}: of the emails given $\post=0.8$, about
$80\%$ should be spam. \code{calibration.py} implements the log-loss, the Brier score, the expected
calibration error (ECE: the average gap between forecast and observed spam rate over 10 bins),
reliability diagrams, temperature scaling, $\post_T=\sigm(\logit\post/T)$, which makes a model more
($T<1$) or less ($T>1$) confident, and Platt scaling, $\post'=\sigm(A\logit\post+B)$, which
re-calibrates a model on validation data. The tests check every score against hand-worked examples and
scikit-learn.

\begin{table}[!htbp]
\centering\footnotesize\setlength{\tabcolsep}{3.5pt}
\input{\resultsdir/c4_scores.tex}
\caption{Forecast quality on test data (mean $\pm$ standard error, 20 splits). Accuracy uses the cut
$0.5$; Cost and Inspected are for the rule with inspection ($\cfp=10$, $\cfn=1$, $\ci=0.3$). Spambase
uses $\log(1+x)$ features. ``+ Platt'': re-calibrated on validation data.}
\label{tab:c4}
\end{table}

On synthetic data, where both models' assumptions hold, they are as good as the true posterior
(Table~\ref{tab:c4}), which checks the code. On Spambase, logistic regression beats the Gaussian
model on every score (log-loss $0.185$ against $0.224$, ECE $0.021$ against $0.039$) and costs
$0.124$ per email against $0.142$, although their accuracies differ by only $1.3$ points. The
reliability diagram (Figure~\ref{fig:c4-rel}) shows where: the Gaussian model underestimates spam for
$\post$ between $0.3$ and $0.6$, exactly where the lower threshold sits. Platt scaling repairs much of
this: the Gaussian model's ECE falls to $0.024$ and its cost to $0.132$.

\begin{figure}[!htbp]
\centering
\includegraphics[width=0.62\linewidth]{c4_reliability.pdf}
\caption{Reliability diagrams on the test parts of all 20 splits. Points above the diagonal mean the
model underestimates spam; the histograms count the forecasts in each bin (log scale).}
\label{fig:c4-rel}
\end{figure}

\paragraph{Same accuracy, different cost.}
To isolate the effect, we took the logistic model and changed only its confidence. Temperature
scaling never moves a probability across $0.5$, so accuracy stays at $93.6\%$, but it changes which
emails pass $0.3$ and $0.97$ (Table~\ref{tab:c4-temp}). Making the model over-confident ($T=0.4$)
raises the number of real emails sent to spam from $3.8$ to $10.5$ per test split and the cost by a
third. Log-loss detects this; the Brier score barely moves.

\begin{table}[!htbp]
\centering\footnotesize\setlength{\tabcolsep}{4pt}
\input{\resultsdir/c4_temperature.tex}
\caption{Logistic regression on Spambase after temperature scaling (mean $\pm$ standard error).}
\label{tab:c4-temp}
\end{table}

So far the test emails came from the same mix as the training emails. A deployed filter cannot count
on that, which led to the next question.
